% =============================================================================
% Chapter 1 --- Normative Source Hierarchy
% =============================================================================
\chapter{Normative Source Hierarchy}
\label{ch:sources}

\section{Why a source hierarchy matters}

Singapore's AI governance landscape comprises multiple overlapping documents
from different agencies (IMDA, PDPC, CSA) with different publication statuses
(official, consultation, reference). Without explicit precedence, an
organisation cannot answer the question \emph{``Which source should I follow
when guidance conflicts?''}

This chapter establishes the authoritative ordering and explains how each
source class contributes to the overall framework.

\section{Source precedence}

Sources are applied in the following order:

\begin{enumerate}
  \item \textbf{Normative Singapore governance baseline}: Official IMDA and
        PDPC guidance on AI governance, agentic AI, LLM testing, and personal
        data use in AI.
  \item \textbf{Normative operational assurance baseline}: Official IMDA/AI
        Verify materials on LLM testing and AI assurance.
  \item \textbf{Supplementary security baseline}: Official CSA secure-AI
        guidance, with explicit notation where guidance is draft or
        consultation-stage.
  \item \textbf{Supporting reference material}: AI Verify open-source
        artefacts and academic references.
\end{enumerate}

\section{Normative Singapore sources}

\begin{table}[htbp]
\centering
\footnotesize
\caption{Normative and Supplementary Source Set}
\label{tab:sources}
\begin{tabularx}{\textwidth}{@{}l l l X@{}}
\toprule
\textbf{ID} & \textbf{Source} & \textbf{Date} & \textbf{Role} \\
\midrule
\sourceid{SG-01} & IMDA MGF for Agentic AI & Jan 2026 & Primary normative baseline for agentic use-case selection, bounded autonomy, human accountability \\
\sourceid{SG-02} & IMDA Starter Kit for LLM Testing & Jan 2026 & Primary baseline for pre-deployment testing structure \\
\sourceid{SG-03} & PDPC Model AI Governance Framework & Jan 2020 & Baseline governance structure \\
\sourceid{SG-04} & PDPC Advisory Guidelines (Personal Data in AI) & Apr 2024 & PDPA-aligned data handling \\
\sourceid{SG-05} & IMDA Project Moonshot / AI Verify & May 2024 & Benchmarking and red-teaming \\
\sourceid{SG-06} & Singapore's Approach to AI Governance & Ongoing & Practical risk-based position \\
\sourceid{SG-07} & CSA Guidelines on Securing AI Systems & Oct 2024 & Supplementary secure-by-design baseline \\
\sourceid{SG-08} & CSA Agentic AI Addendum (Consultation) & Oct 2025 & Supplementary security input (consultation-stage) \\
\bottomrule
\end{tabularx}
\end{table}

\section{Source class descriptions}

\subsection{SG-01: IMDA Model Governance Framework for Agentic AI}

The MGF for Agentic AI is the primary normative source. It defines four
pillars:

\begin{enumerate}
  \item \textbf{Pillar 1}: Assess and Bound Risks Upfront
  \item \textbf{Pillar 2}: Make Humans Meaningfully Accountable
  \item \textbf{Pillar 3}: Implement Technical Controls and Processes
  \item \textbf{Pillar 4}: Enable End-User Responsibility
\end{enumerate}

Each requirement in this specification traces back to one or more MGF sections.

\subsection{SG-02: IMDA Starter Kit for Testing LLM-Based Applications}

Provides the structure for pre-deployment safety testing:

\begin{itemize}
  \item Baseline safety evaluation methodology
  \item Statistical sufficiency requirements
  \item Red-teaming and adversarial testing guidance
\end{itemize}

\subsection{SG-03: PDPC Model AI Governance Framework}

The foundational governance document establishing:

\begin{itemize}
  \item Named accountable ownership requirement
  \item Risk-based approach to AI deployment
  \item Internal governance structures
\end{itemize}

\subsection{SG-04: PDPC Advisory Guidelines on Personal Data in AI}

PDPA-aligned guidance for personal data handling:

\begin{itemize}
  \item Lawful basis for AI processing
  \item Data minimisation requirements
  \item Retention and deletion obligations
\end{itemize}

\subsection{SG-05: IMDA Project Moonshot / AI Verify}

Operationalisation tooling:

\begin{itemize}
  \item AI Verify benchmarking framework
  \item Moonshot red-teaming capabilities
  \item Testing artefact structure
\end{itemize}

\subsection{SG-07: CSA Guidelines on Securing AI Systems}

Security baseline covering:

\begin{itemize}
  \item Secure-by-design principles
  \item Adversarial resilience
  \item Incident response
\end{itemize}

\subsection{SG-08: CSA Agentic AI Addendum (Consultation)}

\begin{tcolorbox}[colback=sgamber!10,colframe=sgamber!70,title=\faIcon{exclamation-triangle} Consultation-Stage Document]
This source is in public consultation as of October 2025. Requirements derived
from SG-08 are marked as \textbf{supplementary} and may change when the final
version is published.
\end{tcolorbox}

Covers agentic-specific security concerns:

\begin{itemize}
  \item Tool access sandboxing
  \item Agent identity attestation
  \item Least-privilege tool boundaries
\end{itemize}

\section{Source document locations}

\begin{table}[htbp]
\centering
\footnotesize
\caption{Source Document Repository Paths}
\begin{tabularx}{\textwidth}{@{}l X@{}}
\toprule
\textbf{Source ID} & \textbf{Repository Path} \\
\midrule
\sourceid{SG-01} & \filepath{docs/regulations/mgf-for-agentic-ai.pdf} \\
\sourceid{SG-02} & \filepath{docs/regulations/aiverify-main/} \\
\sourceid{SG-05} & \filepath{docs/regulations/moonshot-cicd-main/} \\
\sourceid{SG-06} & \filepath{docs/pdf/technical-reports/06-singapore-ai-safety-framework.pdf} \\
\bottomrule
\end{tabularx}
\end{table}

\section{Cross-referencing peer specifications}

This specification cross-references the following peer documents:

\begin{itemize}
  \item \textbf{Regulations Corpus} (\filepath{docs/latex/specs/regulations/}) --- 
        Agentic AI regulatory requirements and implementation evidence
  \item \textbf{Clinerules-Agents} (\filepath{docs/latex/specs/clinerules-agents/}) --- 
        Agent rule-pack architecture
  \item \textbf{TB-HITL Specification} (\filepath{docs/latex/specs/tb/}) --- 
        Trial Balance human-in-the-loop workflow
\end{itemize}